% =============================================================================
% Halaman hak cipta (kolofon).
%
% CATATAN UNTUK PENYUNTING
% Ketentuan lisensi di bawah menyatakan izin penggunaan untuk keperluan
% pendidikan di lingkungan ITERA. Bila program studi hendak memakai lisensi
% terbuka -- misalnya Creative Commons BY-SA 4.0 -- ganti alinea "Ketentuan
% penggunaan" dan cantumkan lisensinya secara eksplisit. Sesuaikan pula nomor
% ISBN bila dokumen didaftarkan.
% =============================================================================
\thispagestyle{empty}

\vspace*{2.5cm}

\noindent
{\large\bfseries\color{Ink}\jenisdokumen\ \coursename}\par
\vspace{0.25em}
{\color{Slate}\coursecode\ --- \doccetakan, Versi \docversion}\par

\vspace{1.2cm}
\noindent\brandrule
\vspace{1.2cm}

\noindent
\begin{minipage}{\linewidth}
\small

\textbf{Hak Cipta \textcopyright\ 2026}\par
\pemeganghakcipta\par

\vspace{1.0em}

\textbf{Penyusun}\par
\penyusun\par

\vspace{1.0em}

\textbf{Penerbit}\par
Program Studi \programname\par
\facultyname, \universityname\par
Jalan Terusan Ryacudu, Way Huwi, Lampung Selatan 35365\par

\vspace{1.0em}

\textbf{Ketentuan penggunaan}\par
Hak cipta atas dokumen ini dimiliki oleh \pemeganghakcipta.
Dokumen diperbanyak dan digunakan untuk keperluan pembelajaran di lingkungan
\universityname. Penggunaan di luar lingkungan tersebut, termasuk penggandaan
untuk kepentingan komersial, memerlukan izin tertulis dari pemegang hak cipta.

\vspace{0.8em}

Kutipan untuk keperluan pendidikan dan penelitian diperbolehkan sepanjang
menyebutkan sumbernya. Kode program, skrip SQL, dan berkas konfigurasi yang
menyertai dokumen ini tersedia pada repositori yang disebut di bawah dan tunduk
pada ketentuan lisensi yang tercantum di sana.

\vspace{1.0em}

\textbf{Repositori}\par
\url{\repositori}

\vspace{1.0em}

\textbf{Dokumen pendamping}\par
\bukupendamping

\vspace{1.0em}

\textbf{Cara mengutip}\par
\penyusun. (2026). \emph{\jenisdokumen\ \coursename} (Versi \docversion).
Program Studi \programname, \universityname.

\vspace{1.2em}


\end{minipage}

\vfill
\noindent\brandrule
\clearpage
